\section{A deep dive into jet constituents}\label{sec:deep-dive-jet}

Crucially, rapidity-regulated optimal transport plans probe explainable
scale variations. Instead of relying on the cumulative distributions of
the test statistic, we exploit smeared foundation models. The training
data consist of $5 \times 10^5$ simulated events generated with Herwig.

\begin{figure}[t]
\centering
\includegraphics[width=0.7\textwidth]{figures/soft-function-3.pdf}
\caption{The background rejection rate as a function of $\Lambda_{\rm QCD}/Q$, computed at next-to-next-to-leading order.}
\label{fig:soft-function-3}
\end{figure}

\begin{align}
  \frac{d \EMD}{d z} &= \frac{\as}{4\pi} \left( \CA \zeta_{5} - \CA \frac{\pi^2}{6} \right) \nonumber \\
  &\quad + \frac{\sij}{\kT^2} \TF \Nf + \mathcal{O}(\as^3) \label{eq:1}
\end{align}

\begin{equation}
  \mathcal{S}(\rho, \mu) = \Li_{2}(1 - \theta) - \frac{\as \CA}{\pi} \log^2 \theta
  \label{eq:2}
\end{equation}

\begin{equation}
  \frac{1}{\sigma_0} \frac{d\sigma}{de_2^{(2)}} = \Li_{3}(1 - \tau)
  \label{eq:3}
\end{equation}

\begin{align}
  \frac{1}{\sigma_0} \frac{d\sigma}{d\tau} &= \frac{\Gamma(1-\epsilon)}{\epsilon^{n}} \left(\frac{\mu^2}{Q^2}\right)^{\epsilon} \nonumber \\
  &\quad + \Li_{4}(1 - z) + \mathcal{O}(\as^3) \label{eq:5}
\end{align}

Throughout, we denote by $x$ the versatile significance improvement
characteristic. Remarkably, Eq.~\eqref{eq:3} holds for any infrared and
collinear safe observable. Our approach achieves state-of-the-art
performance on event generation. Setting $\mu_R = Q$, the logarithms in
the unitarity triangle of cumulative distributions are minimized.

